% GENERATED FILE. Edit canonical lessons, then run npm run generate:books.
% canonical-source-hash: fnv1a64:a08f6cfb7758f9aa
% canonical-lessons: ML-C459-cumattuka, ML-C459-pirikkuka, ML-C459-utpadippikkuka, ML-C459-pradesika, ML-C459-desiya

\chapter{Taxes and Production}
\label{ch:ml-taxes-and-production}
\hlchaptermodality{459}

\begin{chapteropening}
\textbf{By the end of this chapter:} \emph{I can say to impose, to collect, to produce, local and national.}

Complete the last lesson of chapter 459: I can say to impose, to collect, to produce, local and national.
\end{chapteropening}

\section[\texorpdfstring{\ml{ചുമത്തുക} (cumattuka)}{ചുമത്തുക (cumattuka)}]{\ml{ചുമത്തുക} (cumattuka) — to impose; to levy}

\label{lesson:ML-C459-cumattuka}



\begin{quote}
Before the new one: say the Malayalam for a loan, then the Malayalam for to donate.
\end{quote}

\subsection*{You'll want to know: \ml{ചുമത്തുക}}
\textbf{\ml{ചുമത്തുക}} — \emph{cumattuka} — \textquotedblleft{}to impose; to levy\textquotedblright{}.

\begin{quote}
\ml{സർക്കാർ} \ml{പെട്രോളിന്} \ml{പുതിയ} \ml{നികുതി} \ml{ചുമത്തി}.
\end{quote}

\begin{quote}
\emph{sarkkār peṭrōḷinŭ putiya nikuti cumatti.}
\end{quote}

\begin{quote}
The government imposed a new tax on petrol.
\end{quote}

Past: \ml{ചുമത്തി}, \emph{cumatti}.

Also to bring a charge: \ml{കുറ്റം} \ml{ചുമത്തുക}, to charge with a crime.

\begin{cousinweb}
\textbf{\ml{ചുമത്തുക}} (\emph{cumattuka}) is native Dravidian: DEDR 2677 lists Malayalam \emph{cumattuka}, \textquotedblleft{}to load, to charge\textquotedblright{}, beside \emph{cuma}, \textquotedblleft{}a burden\textquotedblright{}: to load a burden on someone.
\end{cousinweb}

\begin{grammarlens}[title={What you've built}]
One more word for saying what happens if.
\end{grammarlens}

\subsection*{Your turn}
\begin{itemize}
\raggedright
  \item \emph{Say it:} \emph{cumattuka}
  \item \emph{Say it:} \emph{cumattuka}, once more
  \item \emph{Say it:} \emph{sarkkār peṭrōḷinŭ putiya nikuti cumatti.} — The government imposed a new tax on petrol.
  \item \emph{Recall:} say the Malayalam for a loan, then the Malayalam for to donate, then say \emph{cumattuka} again
\end{itemize}

\begin{tcolorbox}[breakable,colback=teal!4,colframe=teal!35!black,title={Before you move on}]
What does \ml{ചുമത്തുക} mean? (\textquotedblleft{}To impose; to levy\textquotedblright{}.) Say it once more.
\end{tcolorbox}
\section[\texorpdfstring{\ml{പിരിക്കുക} (pirikkuka)}{പിരിക്കുക (pirikkuka)}]{\ml{പിരിക്കുക} (pirikkuka) — to collect dues}

\label{lesson:ML-C459-pirikkuka}



\begin{quote}
Before the new one: say the Malayalam for to donate, then the Malayalam for to impose.
\end{quote}

\subsection*{You'll want to know: \ml{പിരിക്കുക}}
\textbf{\ml{പിരിക്കുക}} — \emph{pirikkuka} — \textquotedblleft{}to collect dues\textquotedblright{}.

\begin{quote}
\ml{പഞ്ചായത്ത്} \ml{വീടുകളിൽ} \ml{നിന്ന്} \ml{നികുതി} \ml{പിരിച്ചു}.
\end{quote}

\begin{quote}
\emph{pañcāyattŭ vīṭukaḷil ninnŭ nikuti piriccu.}
\end{quote}

\begin{quote}
The panchayat collected tax from the houses.
\end{quote}

Past: \ml{പിരിച്ചു}, \emph{piriccu}.

It also means to dismiss: \ml{അവനെ} \ml{ജോലിയിൽ} \ml{നിന്ന്} \ml{പിരിച്ചുവിട്ടു}, they dismissed him from work.

\begin{cousinweb}
\textbf{\ml{പിരിക്കുക}} (\emph{pirikkuka}) is native Dravidian: DEDR 4176 lists Malayalam \emph{pirikka}, \textquotedblleft{}to sever, to dismiss\textquotedblright{}. DEDR does not list the sense \textquotedblleft{}to collect dues\textquotedblright{}.
\end{cousinweb}

\begin{grammarlens}[title={What you've built}]
One more word for saying what happens if.
\end{grammarlens}

\subsection*{Your turn}
\begin{itemize}
\raggedright
  \item \emph{Say it:} \emph{pirikkuka}
  \item \emph{Say it:} \emph{pirikkuka}, once more
  \item \emph{Say it:} \emph{pañcāyattŭ vīṭukaḷil ninnŭ nikuti piriccu.} — The panchayat collected tax from the houses.
  \item \emph{Recall:} say the Malayalam for to donate, then the Malayalam for to impose, then say \emph{pirikkuka} again
\end{itemize}

\begin{tcolorbox}[breakable,colback=teal!4,colframe=teal!35!black,title={Before you move on}]
What does \ml{പിരിക്കുക} mean? (\textquotedblleft{}To collect dues\textquotedblright{}.) Say it once more.
\end{tcolorbox}
\section[\texorpdfstring{\ml{ഉത്പാദിപ്പിക്കുക} (utpādippikkuka)}{ഉത്പാദിപ്പിക്കുക (utpādippikkuka)}]{\ml{ഉത്പാദിപ്പിക്കുക} (utpādippikkuka) — to manufacture; to produce}

\label{lesson:ML-C459-utpadippikkuka}



\begin{quote}
Before the new one: say the Malayalam for to impose, then the Malayalam for to collect dues.
\end{quote}

\subsection*{You'll want to know: \ml{ഉത്പാദിപ്പിക്കുക}}
\textbf{\ml{ഉത്പാദിപ്പിക്കുക}} — \emph{utpādippikkuka} — \textquotedblleft{}to manufacture; to produce\textquotedblright{}.

\begin{quote}
\ml{ആ} \ml{ഫാക്ടറി} \ml{ദിവസവും} \ml{ആയിരം} \ml{സൈക്കിളുകൾ} \ml{ഉത്പാദിപ്പിക്കുന്നു}.
\end{quote}

\begin{quote}
\emph{ā phākṭaṟi divasavuṁ āyiraṁ saikkiḷukaḷ utpādippikkunnu.}
\end{quote}

\begin{quote}
That factory produces a thousand bicycles every day.
\end{quote}

Past: \ml{ഉത്പാദിപ്പിച്ചു}, \emph{utpādippiccu}.

The noun is \ml{ഉത്പാദനം}, production.

\begin{cousinweb}
\textbf{\ml{ഉത്പാദിപ്പിക്കുക}} (\emph{utpādippikkuka}) is built from Sanskrit \emph{utpādana}, \textquotedblleft{}bringing forth, producing\textquotedblright{} (Monier-Williams), with the Malayalam verb ending -\ml{ിപ്പിക്കുക}.
\end{cousinweb}

\begin{grammarlens}[title={What you've built}]
One more word for saying what happens if.
\end{grammarlens}

\subsection*{Your turn}
\begin{itemize}
\raggedright
  \item \emph{Say it:} \emph{utpādippikkuka}
  \item \emph{Say it:} \emph{utpādippikkuka}, once more
  \item \emph{Say it:} \emph{ā phākṭaṟi divasavuṁ āyiraṁ saikkiḷukaḷ utpādippikkunnu.} — That factory produces a thousand bicycles every day.
  \item \emph{Recall:} say the Malayalam for to impose, then the Malayalam for to collect dues, then say \emph{utpādippikkuka} again
\end{itemize}

\begin{tcolorbox}[breakable,colback=teal!4,colframe=teal!35!black,title={Before you move on}]
What does \ml{ഉത്പാദിപ്പിക്കുക} mean? (\textquotedblleft{}To manufacture; to produce\textquotedblright{}.) Say it once more.
\end{tcolorbox}
\section[\texorpdfstring{\ml{പ്രാദേശിക} (prādēśika)}{പ്രാദേശിക (prādēśika)}]{\ml{പ്രാദേശിക} (prādēśika) — local; regional}

\label{lesson:ML-C459-pradesika}



\begin{quote}
Before the new one: say the Malayalam for to collect dues, then the Malayalam for to manufacture.
\end{quote}

\subsection*{You'll want to know: \ml{പ്രാദേശിക}}
\textbf{\ml{പ്രാദേശിക}} — \emph{prādēśika} — \textquotedblleft{}local; regional\textquotedblright{}.

\begin{quote}
\ml{അവൻ} \ml{ഒരു} \ml{പ്രാദേശിക} \ml{പത്രത്തിൽ} \ml{ജോലി} \ml{ചെയ്യുന്നു}.
\end{quote}

\begin{quote}
\emph{avan oru prādēśika patrattil jōli ceyyunnu.}
\end{quote}

\begin{quote}
He works for a local newspaper.
\end{quote}

It stands before a noun: \ml{പ്രാദേശിക} \ml{ഭാഷ}, a regional language.

\begin{cousinweb}
\textbf{\ml{പ്രാദേശിക}} (\emph{prādēśika}) is borrowed from Sanskrit \emph{prādeśika}, \textquotedblleft{}local, limited\textquotedblright{} (Monier-Williams), made from \emph{pradeśa}, \textquotedblleft{}a spot, a region\textquotedblright{}.
\end{cousinweb}

\begin{grammarlens}[title={What you've built}]
One more word for saying what happens if.
\end{grammarlens}

\subsection*{Your turn}
\begin{itemize}
\raggedright
  \item \emph{Say it:} \emph{prādēśika}
  \item \emph{Say it:} \emph{prādēśika}, once more
  \item \emph{Say it:} \emph{avan oru prādēśika patrattil jōli ceyyunnu.} — He works for a local newspaper.
  \item \emph{Recall:} say the Malayalam for to collect dues, then the Malayalam for to manufacture, then say \emph{prādēśika} again
\end{itemize}

\begin{tcolorbox}[breakable,colback=teal!4,colframe=teal!35!black,title={Before you move on}]
What does \ml{പ്രാദേശിക} mean? (\textquotedblleft{}Local; regional\textquotedblright{}.) Say it once more.
\end{tcolorbox}
\section[\texorpdfstring{\ml{ദേശീയ} (dēśīya)}{ദേശീയ (dēśīya)}]{\ml{ദേശീയ} (dēśīya) — national}

\label{lesson:ML-C459-desiya}



\begin{quote}
Before the new one: say the Malayalam for to manufacture, then the Malayalam for local.
\end{quote}

\subsection*{You'll want to know: \ml{ദേശീയ}}
\textbf{\ml{ദേശീയ}} — \emph{dēśīya} — \textquotedblleft{}national\textquotedblright{}.

\begin{quote}
\ml{ഇന്ന്} \ml{ഒരു} \ml{ദേശീയ} \ml{അവധി} \ml{ദിവസമാണ്}.
\end{quote}

\begin{quote}
\emph{innŭ oru dēśīya avadhi divasamāṇŭ.}
\end{quote}

\begin{quote}
Today is a national holiday.
\end{quote}

It stands before a noun: \ml{ദേശീയ} \ml{പതാക}, the national flag.

\begin{cousinweb}
\textbf{\ml{ദേശീയ}} (\emph{dēśīya}) is borrowed from Sanskrit \emph{deśīya}, \textquotedblleft{}belonging to a country, native\textquotedblright{} (Monier-Williams), from \emph{deśa}, \textquotedblleft{}a country\textquotedblright{}.
\end{cousinweb}

\begin{grammarlens}[title={What you've built}]
One more word for saying what happens if.
\end{grammarlens}

\subsection*{Your turn}
\begin{itemize}
\raggedright
  \item \emph{Say it:} \emph{dēśīya}
  \item \emph{Say it:} \emph{dēśīya}, once more
  \item \emph{Say it:} \emph{innŭ oru dēśīya avadhi divasamāṇŭ.} — Today is a national holiday.
  \item \emph{Recall:} say the Malayalam for to manufacture, then the Malayalam for local, then say \emph{dēśīya} again
\end{itemize}

\begin{tcolorbox}[breakable,colback=teal!4,colframe=teal!35!black,title={Before you move on}]
What does \ml{ദേശീയ} mean? (\textquotedblleft{}National\textquotedblright{}.) Say it once more.
\end{tcolorbox}
